\documentclass[9pt,conference]{IEEEtran}

\usepackage[margin=0.4in]{geometry}
\usepackage{amsmath,amssymb}
\usepackage{booktabs}
\usepackage{array}
\usepackage{graphicx}
\usepackage{caption}
\usepackage{xcolor}
\usepackage[colorlinks=true,linkcolor=blue!55!black,citecolor=green!40!black,urlcolor=blue!60!black]{hyperref}
\usepackage{cleveref}
\usepackage{enumitem}
\setlist{noitemsep,topsep=1pt,parsep=0pt,partopsep=0pt,leftmargin=*}

\begin{document}

\title{Two Hypotheses for Agentic RAG: Trust Calibration and Learning Outcomes}

\author{
  \IEEEauthorblockN{Coen Petto \and Cameron Suess}
}

\maketitle

\begin{abstract}
One system manipulation, two experiments. System A escalates to a
human at seeded points of low-confidence or contradictory retrieval,
its policy modeled on a minimal-proactivity design pattern; System B
never escalates.
We state two hypotheses built on this manipulation: H1, that experts
grading system output calibrate their trust better, and choose more
defensibly, under System A; and H2, that novices learn more from
System A. Every clause is grounded and cited by author and year, with
a side-by-side table of the two drafted experiments, a verified
quote sheet, and a September through December execution plan.
\end{abstract}

\section{Two Hypotheses}
\label{sec:hyp}

Both experiments share one manipulation on the same proprietary
corpus. System A's escalation policy is modeled on, not a full
implementation of, Minimal Proactivity Policies (Sánchez-Adame and
Mendoza, 2026): initiations should be bounded and brief, immediately
reversible, carry an explicit provenance cue naming the source that
justified the pause, and default to silence rather than speculation
when evidence is insufficient. System B always answers from its top
retrieval, with no pause and no provenance cue, matching the anchor
paper's non-escalating baseline (Mishra et al., 2026). Because
agentic reasoning's fluency can inflate trust independent of actual
accuracy (Möhle et al., 2026), grading and choice, not stated
confidence, are what Experiment 1 measures. These three papers are
sources and design guides, not specifications this project is bound
to reproduce in full; the three-month build-and-run window means each
is adapted to what two people can actually implement and run.

\begin{quote}\small
\textbf{H1 (Trust).} Among domain experts who grade individual
retrieval-and-response outputs from System A and System B on a shared
proprietary corpus, (a) experts' grades will track true output
accuracy more closely for System A than System B, and (b) when asked
at the end of the session which system they would choose for an
upcoming task, and why, experts will choose System A more often and
justify the choice by citing verifiable evidence (flagged uncertainty,
cited passages) rather than confidence or fluency alone.
\end{quote}

\begin{itemize}
\item \textbf{``Grade individual outputs from System A and System
  B''}: directly modeled on a published human-in-the-loop grading
  study where participants ``assumed the role of the external
  commander and evaluated plans provided by the internal robot \ldots
  and rate those plans as optimal or suboptimal'' (Chakraborti,
  Sreedharan, Grover, and Kambhampati, 2018).
\item \textbf{``Track true output accuracy''}: this is trust
  calibration, defined as ``the alignment between users' trust and
  the system's underlying capability for a given task'' (Möhle et
  al., 2026), the same quantity named as an evaluation criterion,
  Expected Calibration Error, for the open problem this project
  addresses (Mishra et al., 2026).
\item \textbf{``Citing verifiable evidence rather than confidence or
  fluency alone''}: Möhle et al. (2026) find a significant
  interaction where stated accuracy shifts trust more strongly when
  the accompanying reasoning is convincing, evidence that convincing
  reasoning is not a generic trust booster but a conditioning cue on
  how capability claims become trust. System A's provenance
  requirement (Sánchez-Adame and Mendoza, 2026) gives experts exactly
  the verifiable evidence that fluency alone does not.
\item \textbf{Choice plus justification, not a scale alone}: extends
  the end-of-study items of Chakraborti et al. (2018), ``I trust the
  robot to work on its own'' and ``my trust in the robot increased
  during the study,'' into a forced choice with a coded free-text
  reason, analyzed by the course's thematic-analysis method (CS 567
  slides, Ortega).
\end{itemize}

\newpage
\begin{quote}\small
\textbf{H2 (Learning).} Among novice engineers completing a short
onboarding checklist for the same proprietary tool with System A or
System B, System A will produce a higher score on an unassisted,
task-embedded next-step post-test than System B.
\end{quote}

\begin{itemize}
\item \textbf{``Novice engineers, onboarding checklist, proprietary
  tool''}: proprietary information is the actual RAG use case and the
  source of miss-retrieval risk (the project's RAG guide), matching
  the mission this project was set (the author's project notes).
  Section~\ref{sec:corpus} gives the corpus construction method;
  Section~\ref{sec:env} gives the environment and the outcome measure.
\item \textbf{Agentic manipulation}: System A versus System B
  satisfies the field's own ``chatbot test'': a study whose result
  would be unchanged with the agency stripped out has no independent
  variable (CS 567 slides, Ortega).
\end{itemize}

\begin{table*}[t]
\centering\small
\begin{tabular}{@{}p{0.10\linewidth}p{0.43\linewidth}p{0.43\linewidth}@{}}
\toprule
 & \textbf{Experiment 1: Trust (H1)} & \textbf{Experiment 2: Learning (H2)} \\
\midrule
Population & Domain experts grading system output & Novice engineers (new hires) onboarding \\
Task & Grade a fixed set of retrieval-and-response outputs, half from System A and half from System B & Complete a short onboarding checklist for the invented tool, using System A or System B \\
Procedure & Blind, randomized-order grading of outputs fixed in advance (Chakraborti et al., 2018; Möhle et al., 2026); end-of-session forced choice of A or B for an upcoming task, with a free-text reason & Pull-based query loop (Section~\ref{sec:env}); two or three checklist items seeded with low-confidence or contradictory retrieval \\
Primary measure & Grade-accuracy correlation per system (calibration) & Unassisted, task-embedded next-step post-test accuracy (Section~\ref{sec:env}) \\
Secondary measures & Final choice (A/B); thematic code of the justification (evidence-based vs.\ confidence/fluency-based); Competence/Benevolence/Integrity trust scale (Möhle et al., 2026) & Completion time; NASA-TLX (optional) \\
Design & Within-subjects, counterbalanced order, $N \approx 15$--20 & Within-subjects, two-level independent variable, $N \approx 15$--20 \\
Session length & $\sim$15--20 min; no live model prototype needed, outputs are pre-recorded & $\sim$30--45 min \\
\bottomrule
\end{tabular}
\caption{The two drafted experiments, sharing System A (escalates, its policy modeled on MiPP) versus System B (never escalates) on the same proprietary corpus.}
\label{tab:experiments}
\end{table*}

\section{Experiment 2: Environment and Outcome Measure}
\label{sec:env}

The participant is framed as a new hire completing a short onboarding
checklist for an invented internal tool, using the RAG assistant on
their own initiative; the assistant never contacts the participant
first. In the System A condition, seeded low-confidence or
contradictory items trigger a pause and a request to confirm before
proceeding, shown in Fig.~\ref{fig:hitl}; in the System B condition,
the same items are answered directly, without a flag.

The outcome measure is task-based, not a declarative quiz: a short,
unassisted, procedural post-test at the end of the session, given a
checklist state the participant has not seen answered, what is the
correct next step. This follows Ortega's own finding that
self-reported learning and objective task performance dissociate:
richer training materials inflated participants' confidence that they
would remember correctly with no change in their actual ability to
identify the next step (Clegg et al., 2022), exactly the illusion
System B is predicted to produce. How the assistant should present or
time an escalation once triggered is left as an open design problem
this initial study motivates.

\begin{figure}[h]
\centering
\includegraphics[width=\linewidth]{fig5_sok_hitl.png}
\caption{The Human-in-the-Loop escalation loop, taken directly from
Mishra et al. (2026, their Fig.~5). System A's policy for this
escalation branch is modeled on MiPP (Sánchez-Adame and Mendoza,
2026); System B disables the branch and always proceeds to Final
Response.}
\label{fig:hitl}
\end{figure}

\section{Corpus Construction and Contamination Check}
\label{sec:corpus}
\small
Following Kirchenbauer et al. (2026), the invented tool used by both
experiments is built in five stages: a short seed premise for the
tool and its purpose; a fictsheet enumerating its commands, flags,
and error conditions; multiple document renderings (a README, a
changelog, an internal wiki page, a deprecated guide left in place to
seed contradiction); question and answer pairs with span-constrained
answers for Experiment 2's post-test and Experiment 1's grading set;
and a blind-versus-informed annotation pass, where the pipeline's own
generator model is asked each item with no supporting document in
context, then again with the document supplied. Items the model
answers correctly blind are discarded as contaminated or guessable;
items it fails blind and passes informed confirm the corpus is the
sole source of the fact, satisfying both the corpus-construction and
the no-prior-knowledge requirements with one procedure.

\section{Verified Reference Sheet}
\small
Every quote below was read directly from the cited primary source,
not taken from a secondary citation.

\begin{itemize}
\item \textbf{Mishra et al. (2026), HITL pattern:} ``This pattern
  models human oversight as a callable API within the action space.
  When epistemic uncertainty exceeds a defined threshold, the policy
  pauses execution to request disambiguation or supervision.''
\item \textbf{Möhle et al. (2026), trust calibration:} ``Trust
  calibration denotes the alignment between users' trust and the
  system's underlying capability for a given task.''
\item \textbf{Möhle et al. (2026), key finding:} ``Results show a
  positive direct effect of stated accuracy on trust and a
  significant interaction: accuracy shifts trust more strongly when
  reasoning is convincing.''
\item \textbf{Sánchez-Adame and Mendoza (2026), MiPP contract, C3:}
  ``Every initiation must expose an explicit provenance cue
  indicating the observable source(s) that justified the initiation
  \ldots enough local information for the user to judge whether the
  initiation is legitimate in context.''
\item \textbf{Sánchez-Adame and Mendoza (2026), MiPP contract, C5:}
  ``When required inputs are missing or ambiguous, the default is
  non-initiation \ldots rather than `trying anyway' and risking
  speculative interruptions, the system should defer \ldots or remain
  silent.''
\item \textbf{Chakraborti, Sreedharan, Grover, and Kambhampati
  (2018), Study 2 design:} ``participants assumed the role of the
  external commander and evaluated plans provided by the internal
  robot \ldots and rate those plans as optimal or suboptimal based on
  that explanation.''
\item \textbf{Chakraborti et al. (2018), end-of-study trust items:}
  ``I trust the robot to work on its own.'' ``My trust in the robot
  increased during the study.''
\item \textbf{CS 567 slides (Ortega), ``the chatbot test'':} ``If
  your study would come out the same, you are studying a chatbot. The
  agentic part is decoration, and the proposal has no independent
  variable.''
\item \textbf{Clegg et al. (2022), general discussion:}
  ``participants rated the chances they would remember how to
  assemble objects as greater than 2/3 of the time on average,
  whereas their actual ability showed that only about 1/3 of the time
  were they able to correctly completely identify all the critical
  information.''
\item \textbf{Kirchenbauer et al. (2026), verifying no prior
  knowledge:} ``prompting the same model used in the data generation
  process to answer the questions in two ways: blind with only the
  question in context, and informed via in-context access to the
  fictional document \ldots we try to ensure that the questions are
  not answerable by a powerful language model that has never even
  seen the fictional documents.''
\item \textbf{The author's project notes:} ``Min Docket: Trust in RAG
  systems. Efficiency in learning with RAG systems.''
\end{itemize}

\section{Execution Plan (September through December)}
\small
\begin{tabular}{@{}p{0.11\linewidth}p{0.85\linewidth}@{}}
\toprule
\textbf{Month} & \textbf{Deliverable} \\
\midrule
Sept, wks 1-2 & File CSU IRB protocol covering both experiments, since
  approval precedes recruitment (CS 567 slides); finalize System A's
  escalation policy and its signal thresholds (Section~\ref{sec:hyp}); read
  Mishra et al. (2026), Möhle et al. (2026), and Sánchez-Adame and
  Mendoza (2026) in full as a team. \\
Sept, wks 3-4 & Build the invented tool and its documents following the
  FictionalQA pipeline (Section~\ref{sec:corpus}); run the
  blind-versus-informed contamination check; write the fixed System
  A/System B response set shared by both experiments so response
  quality is controlled, not sampled (the project's RAG guide). \\
Oct, wks 1-2 & Build Experiment 1's grading interface (static outputs,
  no live model call) and Experiment 2's scripted two-arm prototype;
  instrument logging before piloting. \\
Oct, wks 3-4 & Submit final IRB documents and confirm approval; pilot
  both experiments with 2-3 participants each and fix timing, logging,
  and instructions. \\
Nov, wks 1-3 & Run Experiment 1 (target $N \approx 15$--20 experts) and
  Experiment 2 (target $N \approx 15$--20 novices) in parallel,
  counterbalanced order. \\
Nov, wk 4 & Close data collection; compute grade-accuracy correlation
  and choice/justification codes for Experiment 1, and post-test
  accuracy for Experiment 2; check normality for both. \\
Dec, wks 1-2 & Paired-comparison analysis for each experiment (Wilcoxon
  signed-rank or paired $t$-test depending on the normality check);
  effect size and confidence interval; thematic coding of Experiment
  1's justifications. \\
Dec, wks 3-4 & Write the Method and Results sections of the paper
  draft for both experiments; internal review; scope the
  adaptive-escalation follow-up named as the anchor paper's own next
  step. \\
\bottomrule
\end{tabular}

\enlargethispage{\baselineskip}
\footnotesize
\textbf{References} (alphabetical by author; cited above by author and year):
\begin{itemize}
\item Chakraborti, T., Sreedharan, S., Grover, S., \& Kambhampati, S. (2018). Plan explanations as model reconciliation: An empirical study. arXiv:1802.01013.
\item Clegg, B. A., Karduna, A., Holen, E., Garcia, J., Rhodes, M. G., \& Ortega, F. R. (2022). Multimedia and immersive training materials influence impressions of learning but not learning outcomes. I/ITSEC 2022.
\item CS 567 slides, Ortega, F. R. (2026). Research Problems and Experiments. CS 567 course slides, Colorado State University.
\item Kirchenbauer, J., Mongkolsupawan, N., Wen, Y., Goldstein, T., \& Ippolito, D. (2026). FictionalQA: A dataset for studying memorization and knowledge acquisition. arXiv:2506.05639.
\item Mishra, S., Niroula, S., Yadav, U., Thakur, D., Gyawali, S., \& Gaire, S. (2026). SoK: Agentic retrieval-augmented generation (RAG): Taxonomy, architectures, evaluation, and research directions. arXiv:2603.07379.
\item Möhle, M. H., Zall, M., Thorwart-Gumpert, T., Elbert, N., Thiesse, F., \& Flath, C. M. (2026). ``Trust me, Bro!'' On the influence of agentic support on trust calibration in AI-assisted decision making. HCII 2026, pp. 472-488.
\item Sánchez-Adame, L. M., \& Mendoza, S. (2026). MiPP: Minimal proactivity policies for timely, low-intrusion proactivity. HCII 2026, pp. 363-382.
\item The project's RAG guide (2026). Evaluating human learning from RAG systems: A literature orientation and design guide. Internal project document.
\item The author's project notes: Petto, C. (2026). Handwritten project notes.
\end{itemize}

\end{document}
